% draft/SMALL_EXPERIMENT.tex
% Small-scale controlled experiment (2 facilities, 22 beds)
% Run: PYTHONPATH=. .venv/bin/python experiments/small_experiment.py
% 30 replications | 730-day horizon | 90-day warm-up
% θ_CommHigh recalibrated: 0.75%/month (133-month mean patience)

\section{Experiments: Small-Scale Setting}
\label{sec:small}

\subsection{Setup}

We construct a two-facility, 22-bed instance to verify the model mechanics at a
scale that is transparent and hand-verifiable before applying VCH-calibrated parameters.
All parameters are fixed; only the assignment policy varies.

\begin{table}[H]
\centering
\caption{Small-scale experiment configuration.}
\label{tab:small_config}
\begin{tabular}{@{}p{3.6cm}p{3.0cm}p{3.6cm}p{3.0cm}@{}}
\toprule
\textbf{Parameter} & \textbf{Value} & \textbf{Parameter} & \textbf{Value} \\
\midrule
Facility A (Cedar Grove)
  & 10 beds, $u_n = -0.10$
  & Facility B (Maple House)
  & 12 beds, $u_n = +0.15$ \\[2pt]
Arrival rate $\lambda$
  & 3 clients/month
  & Vacancy rate $\mu$
  & 2 beds/month \\[2pt]
Overload $(\lambda{-}\mu)/\lambda$
  & 33\%
  & Initial queue $Q_0$
  & 10 clients \\[2pt]
$\theta$ (Acute Care)
  & 50\%/month
  & $\theta$ (Comm.\ Emergency)
  & 30\%/month \\[2pt]
$\theta$ (Comm.\ High)
  & 0.75\%/month
  & $\theta$ (Transfer/Site)
  & 15\%/month \\[2pt]
Simulation horizon
  & 730 days
  & Warm-up period
  & 90 days \\[2pt]
Replications
  & 30
  & Offer turnaround
  & 2 days/decline \\
\bottomrule
\end{tabular}
\end{table}

Cedar Grove has a below-average baseline deferral propensity ($u_n = -0.10$).
Maple House has an above-average propensity ($u_n = +0.15$), making it more
selective for high-intensity clients. No gender restrictions apply; gender
compatibility is handled prior to the offer stage and is not part of the
assignment constraints modelled here.

The facility-specific effect $u_n$ and the client intensity coefficient
$\beta_{r_p} = +0.9$ together determine how much a facility's selectiveness
interacts with client complexity. All logistic coefficients
($\alpha$, $\boldsymbol{\beta}$, $\sigma^2$) are placeholder values estimated
from synthetic data mirroring the VCH PARIS schema. Real values would be
obtained from logistic regression on historical offer outcome records; the
framework is designed for direct coefficient substitution once that data is
available.

The Community High abandonment rate $\theta = 0.75\%$/month (mean patience
$\approx 133$ months) is calibrated so that the Cyclic FIFO policy reproduces
the OSA-reported mean wait of 473 days for non-urgent community clients at
VCH scale (Section~\ref{sec:data}). The small-scale experiment uses the same
rate for consistency.

We compare three policies:

\begin{itemize}
  \item \textbf{FCFS (No Declines):} all facility offers are accepted
        ($\hat{p}_{pn} = 0$). Clients are served first-in-first-out across all
        tiers. This eliminates all bed-loss days and serves as a lower bound on
        wasted capacity.
  \item \textbf{Cyclic FIFO:} vacancies rotate through the four priority tiers
        in equal 25\% rotation, serving the longest-waiting eligible client
        within each tier. This reflects VCH's current administrative policy.
  \item \textbf{Optimised:} at each vacancy, the eligible client minimising
        $\hat{p}_{pn}$ is selected, subject to a fairness window of 20 recent
        placements ($\pm 10\%$ tolerance per tier).
\end{itemize}

\subsection{Results}

\subsubsection*{Overall System Performance}

Table~\ref{tab:small_overall} reports system-wide outcomes averaged over 30 replications.

\begin{table}[H]
\centering
\caption{Overall results, small-scale experiment (30 reps, 730-day horizon, 90-day warm-up).
         Values are mean $\pm$ 95\% CI.}
\label{tab:small_overall}
\begin{tabular}{lrrr}
\toprule
\textbf{Metric} & \textbf{FCFS (No Declines)} & \textbf{Cyclic FIFO} & \textbf{Optimised} \\
\midrule
Placements              & $39 \pm 2$    & $39 \pm 2$    & $40 \pm 2$ \\
Declinations            & $0 \pm 0$     & $8 \pm 1$     & $9 \pm 1$  \\
Abandonments            & $26 \pm 2$    & $25 \pm 2$    & $21 \pm 2$ \\
Bed-Loss Days           & $0.0$         & $16.5 \pm 2.5$ & $17.1 \pm 2.2$ \\
Bed-Loss Days/placement & $0.000$       & $0.420 \pm 0.060$ & $0.430 \pm 0.050$ \\
Declination rate        & $0.000$       & $0.170 \pm 0.020$ & $0.170 \pm 0.020$ \\
Mean wait (days)        & $66.4 \pm 11.4$ & $83.4 \pm 12.4$ & $57.9 \pm 6.8$ \\
Mean queue length       & $5.3 \pm 0.8$ & $6.8 \pm 1.0$ & $7.9 \pm 1.2$ \\
\bottomrule
\end{tabular}
\end{table}

At this scale, the declination rates of Cyclic FIFO and Optimised are
statistically indistinguishable (both $0.170 \pm 0.020$; confidence intervals
fully overlap). The primary gain from the optimised policy is in wait times:
mean wait drops from 83.4 to 57.9 days ($-30.6\%$), driven by large reductions
in the urgent tiers.

\subsubsection*{Per-Tier Results}

\begin{table}[H]
\centering
\caption{Per-tier declination rate, small-scale experiment (mean $\pm$ 95\% CI).}
\label{tab:small_tier_declin}
\begin{tabular}{lrrr}
\toprule
\textbf{Tier} & \textbf{FCFS (No Declines)} & \textbf{Cyclic FIFO} & \textbf{Optimised} \\
\midrule
Transfer/Site Specific & $0.000$ & $0.230 \pm 0.050$ & $0.190 \pm 0.050$ \\
Community High         & $0.000$ & $0.150 \pm 0.040$ & $0.220 \pm 0.040$ \\
Community Emergency    & $0.000$ & $0.150 \pm 0.040$ & $0.130 \pm 0.040$ \\
Acute Care             & $0.000$ & $0.130 \pm 0.030$ & $0.130 \pm 0.030$ \\
\bottomrule
\end{tabular}
\end{table}

\begin{table}[H]
\centering
\caption{Per-tier mean wait in days, small-scale experiment (mean $\pm$ 95\% CI).}
\label{tab:small_tier_wait}
\begin{tabular}{lrrr}
\toprule
\textbf{Tier} & \textbf{FCFS (No Declines)} & \textbf{Cyclic FIFO} & \textbf{Optimised} \\
\midrule
Transfer/Site Specific & $67.8 \pm 12.9$ & $66.4 \pm 13.5$ & $67.8 \pm 19.0$ \\
Community High         & $75.1 \pm 11.4$ & $140.8 \pm 30.4$ & $125.2 \pm 21.1$ \\
Community Emergency    & $62.8 \pm 10.9$ & $60.1 \pm 8.9$   & $35.4 \pm 3.8$ \\
Acute Care             & $50.2 \pm 7.4$  & $41.0 \pm 5.2$   & $22.4 \pm 3.1$ \\
\bottomrule
\end{tabular}
\end{table}

\begin{table}[H]
\centering
\caption{Wait-time improvement from Cyclic FIFO to Optimised, small-scale experiment.}
\label{tab:small_improvement}
\begin{tabular}{lr}
\toprule
\textbf{Metric} & \textbf{\% change (Cyclic FIFO $\to$ Optimised)} \\
\midrule
Mean wait (overall)             & $-30.6\%$ \\
\quad Transfer/Site wait        & $+2.1\%$ \\
\quad Community High wait       & $-11.1\%$ \\
\quad Community Emergency wait  & $-41.1\%$ \\
\quad Acute Care wait           & $-45.3\%$ \\
\midrule
Declination rate  & $0.000$ (no significant change) \\
Bed-Loss Days     & $+3.6\%$ (no significant change) \\
\bottomrule
\end{tabular}
\end{table}

\subsection{Discussion}

\paragraph{Why Acute Care wait falls sharply.}
The provider deferral model assigns a priority coefficient $\beta_{q_p} = -0.2$
per ordinal unit. Acute Care clients have $q_p = 3$, reducing their logit by
$0.6$ relative to Transfer clients ($q_p = 0$). This gives them the lowest
predicted $\hat{p}_{pn}$ of any tier, so the optimised policy selects them
most often. Acute Care mean wait falls from 41.0 to 22.4 days ($-45.3\%$), and
Community Emergency wait falls from 60.1 to 35.4 days ($-41.1\%$). This mirrors
real-world practice: urgent clients who are better matches for available beds
are placed faster when the selection rule accounts for match quality.

\paragraph{Why Community High declination rate rises.}
Community High clients have higher clinical complexity on average (higher $r_p$),
giving them elevated $\hat{p}_{pn}$. The optimised policy avoids them in favour
of lower-risk clients. When the fairness window does force a Community High
selection, it picks the best-matched Community High client available; even so,
that client's $\hat{p}_{pn}$ tends to exceed that of a typical Acute Care client.
The result is a higher observed declination rate for Community High (0.220 vs
0.150 under Cyclic FIFO) alongside a shorter mean wait (125.2 vs 140.8 days,
$-11.1\%$). The fairness window limits this to a bounded increase: no tier is
starved of service.

\paragraph{On FCFS (No Declines) mean wait.}
FCFS (No Declines) eliminates all declinations by construction, but it uses pure
first-come, first-served ordering across all tiers. Acute Care clients who arrive
later than Community High clients must wait behind them regardless of urgency.
This yields a higher Acute Care mean wait (50.2 days) than Cyclic FIFO (41.0
days). Separately, Community High wait is lower under FCFS (No Declines) (75.1
days) than Cyclic FIFO (140.8 days) because FCFS does not cap their service
share at 25\%. Zero bed-loss days and short urgent-tier waits are therefore
distinct objectives requiring different selection rules.

\paragraph{Scale and statistical interpretation.}
With 22 beds and 33\% overload, confidence intervals are wide. At this scale,
the optimised policy does not significantly reduce total declinations or
bed-loss days. The key finding is directional: the optimised policy consistently
reduces urgent-tier wait times without starving lower-urgency tiers. At VCH
scale, with 55 facilities and 28 arrivals per month, both the declination
reduction and the wait-time benefits are amplified and statistically robust
(Section~\ref{sec:vchr}).
